\subsection{Two obstructions to a coarsest weak filtration}
\label{subsec:weak-filtration-counterexamples}
A \emph{weak regular equidimensional filtration} is a finite increasing
chain of closed subsets whose nonempty successive differences are
manifolds without boundary, each of a single dimension.  Its layers
need not be connected.  Refinement is refinement of the associated
layer partitions.  These requirements alone do not force a frontier
condition or singular-locus saturation for closed unions.

\begin{proposition}[A compact counterexample in three-dimensional space]
\label{prop:weak-filtration-r3}
There is a compact semialgebraic set $X\subset\mathbb R^3$ with two
closed filtrations having layers of dimensions $0,1,2$, all embedded
smooth manifolds, such that the only common coarsening of their layer
partitions is $\{X\}$.  The set $X$ is not a manifold of a single
dimension.  Consequently the two filtrations have no common weak
regular coarsening, for either embedded smooth or intrinsic PL
regularity.
\end{proposition}
\begin{proof}
Let
\[
 S=\{(x,y,z):x^2+y^2=z^2(1-z),\ 0\le z\le1\},
 \qquad p=(0,0,0),\qquad \delta=\tfrac14,
\]
and define
\[
 c(t)=(t\sqrt{1-t},0,t)\quad(-\delta\le t\le\delta).
\]
Put $J=c([-\delta,0])$, $C=c([0,\delta])$, $q=c(-\delta)$,
$r=c(\delta)$, and $X=S\cup J$.  The set $S$ is closed and bounded:
the equation bounds $x^2+y^2$ when $z\in[0,1]$.  The arcs are
continuous semialgebraic images of compact intervals.  Thus $X$ is
compact and semialgebraic.  The last coordinate of $c$ is $t$,
so $c$ is an injective smooth parametrization with nonzero tangent.
The defining equation shows that $C\subset S$, whereas the negative
part of $J$ is disjoint from $S$.  Hence $J\cap S=\{p\}$.

There is a global smooth chart for the surface away from $p$.
On the open unit disk $D=\{(u,v):u^2+v^2<1\}$ put $a=1-u^2-v^2$.
The polynomial map
\[
 \Phi:D\longrightarrow S\setminus\{p\},\qquad
 \Phi(u,v)=(au,av,a)
\]
is bijective, with smooth inverse $(x,y,z)\mapsto(x/z,y/z)$.
Indeed every point of $S\setminus\{p\}$ has $z>0$, and its
coordinates satisfy $(x/z)^2+(y/z)^2=1-z$.  This also proves that
$\Phi$ is an embedding: differentiating the inverse identity gives
rank two for its derivative.  Thus $S\setminus\{p\}$ is an embedded
smooth surface without boundary, including the cap $(0,0,1)$.

Consider the closed chains
\[
 \{p,q\}=F_0\subset F_1=J\subset F_2=X,
 \qquad
 \{q,r\}=G_0\subset G_1=J\cup C\subset G_2=X.
\]
Their one-dimensional layers are respectively
$c((-\delta,0))$ and $c((-\delta,\delta))$.  Their
two-dimensional layers are $S\setminus\{p\}$ and $S\setminus C$.
The latter is open in $S\setminus\{p\}$, since $C$ is closed in
$S$ and contains $p$.  All the displayed layers are therefore
embedded smooth manifolds without boundary, with dimensions $0,1,2$.
They are semialgebraically locally Euclidean, so the local
triangulation and pointed Hauptvermutung argument in the proof of
Theorem~\ref{result-14} also makes them intrinsically PL regular;
the classical input is the marked compact Hauptvermutung and
triangulation theorem of Shiota~\cite[Theorems~2.1--2.2, pp.~5--6]{Shiota2010}.

Any partition coarser than both layer partitions must put $p,q$ in
one block, by $F_0$, and $q,r$ in one block, by $G_0$.  Since $r$
belongs to $F_2\setminus F_1$, this same block contains all of
$S\setminus\{p\}$.  Since $p$ belongs to $G_1\setminus G_0$, it
also contains $c((-\delta,\delta))$.  With $q$ these sets exhaust
$X$, so the common coarsening has only one block.

At a point in the interior of $J$ a sufficiently small neighborhood
in $X$ is one-dimensional.  At a point in $S$ away from $C$ it is
two-dimensional.  Invariance of dimension
\cite[Theorem~2.26, p.~126]{Hatcher2002Algebraic} therefore prevents
$X$ from being a manifold of a single dimension, and its indiscrete
partition is not weakly regular.
\end{proof}

The same formulas give a closed and bounded semialgebraic example
over every real closed field, with definable smoothness and the
intrinsic finite PL convention.  The chart and its inverse are
defined over that field, and the final contradiction uses local
definable dimension.  Ordinary compactness is asserted here only
over $\mathbb R$.  This three-dimensional construction has a
one-dimensional part; the next example retains pure dimension two.

\begin{proposition}[A pure two-dimensional compact counterexample]
\label{prop:weak-filtration-pure-r4}
There is a compact semialgebraic set $X\subset\mathbb R^4$, of
local dimension two at every point, with two closed filtrations
whose layers are embedded smooth manifolds of dimensions $0,1,2$
and whose only common coarsening is $\{X\}$.  The set $X$ is not
an intrinsic PL manifold or a topological manifold.  The two
filtrations therefore have no common weak regular coarsening.
\end{proposition}
\begin{proof}
Let $N$ be the north pole of $S^2$ and use stereographic coordinates
$(x,y)$ on $S^2\setminus\{N\}$.  Define
\[
 h(x,y)=
 \frac{\sqrt{(x-1)^2+y^2}+\sqrt{(x+1)^2+y^2}-2}
 {1+x^2+y^2},\qquad h(N)=0.
\]
The triangle inequality gives $h\ge0$, with equality in this chart
precisely on $A=[-1,1]\times\{0\}$.  If
$\rho=\sqrt{x^2+y^2}$, the numerator is at most $2\rho$,
because each distance to $(\pm1,0)$ is at most $\rho+1$.
Hence $0\le h\le2\rho/(1+\rho^2)$, which tends to zero at
$N$.  Thus $h$ is continuous and semialgebraic on $S^2$ and is
smooth away from $N$ and the endpoints $a_-,a_+$ of $A$.

In $\mathbb R^3\times\mathbb R$ set
\[
 S_0=\{(z,0):z\in S^2\},\qquad
 S_h=\{(z,h(z)):z\in S^2\},\qquad X=S_0\cup S_h.
\]
Both sheets are compact semialgebraic copies of the sphere, and
their intersection is $A\cup\{N\}$.  Each has local dimension two
at every point, so $X$ has that same local dimension everywhere.
Let $C\subset S_0$ be the great circle given by $y=0$ in the
finite chart, together with $N$.  Consider
\[
 \begin{array}{lll}
 F_0=\{a_-,a_+,N\},& F_1=A\cup\{N\},& F_2=X,\\
 G_0=\{N\},&G_1=C,&G_2=X.
 \end{array}
\]
Write $P_i=F_i\setminus F_{i-1}$ and
$Q_i=G_i\setminus G_{i-1}$, with $F_{-1}=G_{-1}=\varnothing$.
The zero-dimensional layers are finite sets; $P_1$ is the open
arc between the endpoints, and $Q_1=C\setminus\{N\}$ is a smooth
one-manifold.  The layer $P_2$ is the union of the two disjoint
smooth sheets off $A\cup\{N\}$.  The layer $Q_2$ consists of the
two open hemispheres of $S_0\setminus C$ and the smooth graph
$S_h\setminus(A\cup\{N\})$.  At each point of either two-layer,
the sheets are separated by a positive height locally.  Each layer
is consequently an embedded smooth manifold without boundary;
the same local definable-chart argument gives intrinsic PL
regularity.

The bipartite graph of nonempty intersections of the two partitions
contains the edges
\[
 P_0Q_0,\quad P_0Q_1,\quad P_1Q_1,
 \quad P_2Q_1,\quad P_2Q_2.
\]
They connect all six vertices.  Representatives are, in order, the
north pole, an endpoint of $A$, an interior point of $A$, a point
of the base great circle outside $A\cup\{N\}$, and a point of the
base sphere outside $C$.  In a common coarsening, two intersecting
layers must lie in the same block.  Connectivity therefore forces
the sole block $X$.

At an interior point of $A$ the intrinsic germ of $X$ is two disks
glued along a diameter: projection to the sphere identifies each
sheet with a disk and identifies the sheets exactly on that
diameter.  This local topological gluing description is
independent of their ambient tangency.  A link is the graph
$\Theta_4$ with two vertices joined by four edges.  Orienting the
edges in the same direction gives cellular boundary columns
$(-1,1)$.  The three differences $e_2-e_1,e_3-e_1,e_4-e_1$
form an integral basis for its kernel, so
\[
 H_2(X,X\setminus\{x\};\mathbb Z)
 \cong\widetilde H_1(\Theta_4;\mathbb Z)
 \cong\mathbb Z^3.
\]
A two-manifold has local degree-two homology $\mathbb Z$; a
manifold of another dimension has a different local homology
pattern~\cite[p.~126]{Hatcher2002Algebraic}.  Thus $X$ is not a manifold, and its sole common
coarsening is not regular.
\end{proof}

These examples concern the weak filtration convention just stated.
The canonical construction in Theorem~\ref{result-14} requires a
frontier condition and saturation of singular loci of closed unions.
The second example fails the frontier condition for its literal
layer partition: $\overline{P_1}=A$ contains $a_-$ and $a_+$ but
omits $N$ from $P_0=\{a_-,a_+,N\}$.  Thus the stronger
hypotheses have substantive content.  Requiring a coarser chain
to be a subsequence only further restricts the possible coarsenings,
so the negative conclusions remain valid under that convention
as well.
